\documentclass[conference]{IEEEtran}
\IEEEoverridecommandlockouts

\usepackage{cite}
\usepackage{amsmath,amssymb,amsthm}
\usepackage{algorithm}
\usepackage{algpseudocode}
\usepackage{array}
\usepackage{booktabs}
\usepackage{graphicx}
\usepackage{xcolor}
\usepackage{microtype}
\usepackage{subcaption}
\usepackage[hidelinks]{hyperref}
\usepackage[capitalise]{cleveref}

% Theorem environments
\theoremstyle{plain}
\newtheorem{theorem}{Theorem}
\newtheorem{lemma}[theorem]{Lemma}
\newtheorem{proposition}[theorem]{Proposition}
\newtheorem{corollary}[theorem]{Corollary}
\theoremstyle{definition}
\newtheorem{definition}[theorem]{Definition}
\newtheorem{assumption}[theorem]{Assumption}
\theoremstyle{remark}
\newtheorem{remark}[theorem]{Remark}
\newtheorem{observation}[theorem]{Observation}

% Custom macros
\newcommand{\eps}{\varepsilon}
\newcommand{\R}{\mathbb{R}}
\newcommand{\E}{\mathbb{E}}
\newcommand{\Adv}{\mathrm{Adv}}
\newcommand{\NE}{\mathrm{NE}}
\newcommand{\BR}{\mathrm{BR}}
\newcommand{\PoEV}{\mathrm{PoEV}}
\newcommand{\PoA}{\mathrm{PoA}}
\newcommand{\PEG}{\mathrm{PEG}}
\newcommand{\iDP}{i\text{-}\mathrm{DP}}
\newcommand{\DP}{\mathrm{DP}}
\newcommand{\MIA}{\mathrm{MIA}}
\newcommand{\GS}[1]{|G_{#1}|}

% natbib-style citet via cite package fallback:
\providecommand{\citet}[1]{\cite{#1}}

% Render IEEEtran's \paragraph as inline-bold (no a)/b)/c) labels)
\makeatletter
\renewcommand{\paragraph}[1]{\par\medskip\noindent\textbf{#1}\hspace{0.5em}}
\makeatother

% Title and authors
\title{The Privacy Externality Game:\\
Strategic Budget Choices in Sampling-Based Individual Differential Privacy}

\author{
\IEEEauthorblockN{Anonymous Submission}
}

\begin{document}
\maketitle

\begin{abstract}
Individual differential privacy (iDP) gives each data contributor a personal
privacy budget. The currently dominant realisation---sampling-based iDP---turns
each contributor's chosen budget into a per-sample inclusion rate, with a single
noise multiplier shared across the dataset. We show that this realisation
creates a privacy externality: lowering one's own budget mechanically raises
the realised membership-inference (MIA) advantage on every other data point,
even though every individual $(\eps_i, \delta)$-iDP guarantee continues to
hold. We model the resulting strategic situation as a normal-form
\emph{Privacy Externality Game} on a continuous action space, with each
player's utility trading off a shared model-quality term $V(\sigma)$ against
their own realised MIA advantage $A_i(\boldsymbol{\eps})$. We motivate the
shape of $V$ via the excess-empirical-risk bound for convex differentially
private empirical risk minimisation, used as a directional proxy in the
(possibly non-convex) deep-learning regime. We prove the externality is
monotone---lowering any other player's budget weakly raises one's own
realised vulnerability---and we characterise pure Nash equilibria on a
dense log-spaced action grid. \textbf{Our headline finding is that even
under fully symmetric players, the sampling-based iDP game admits no
diagonal Nash equilibrium: at every pure NE one player carries a strictly
higher realised MIA advantage in exchange for a higher shared model
utility, without any coordination between the players.} The
coalition-Stackelberg limit of this asymmetry recovers a recent
budget-manipulation attack on sampling-based iDP as the rational best
response of any coordinated subset of data holders.
\end{abstract}

\begin{IEEEkeywords}
Differential privacy, individual differential privacy, game theory, Nash
equilibrium, collusion, membership inference.
\end{IEEEkeywords}

\input{tex/01_intro}
\input{tex/06_related}
\input{tex/02_background}
\input{tex/03_game}
\input{tex/04_analysis}
\input{tex/05_experiments}
\input{tex/07_discussion}

\bibliographystyle{IEEEtran}
\bibliography{refs}

\input{tex/08_appendix}
\input{tex/09_bayesian}

\end{document}
